\section{Introduction}\label{introduction}

The advent of large reasoning models (LRMs) has redefined the benchmark requiring complex reasoning \cite{guo2025deepseek, jaech2024openai}. 
The efficacy of LRMs has been attributed to the additional computation allocated at test time through explicit reasoning paths, typically encapsulated within delimiters such as \texttt{<think>} ... \texttt{</think>}. By exploring logical traces, reflecting on internal errors, and self-verifying intermediate steps, these models have achieved breakthroughs in mathematical and symbolic reasoning tasks. These days, this \textit{think-before-answer} paradigm has emerged as the de facto standard.

\begin{figure}[t]
    \centering
    \includegraphics[width=\linewidth]{ACL2026/figures/0_landing_img.pdf}
    \caption{Overview of our training pipeline. We generate paired reasoning samples and non-reasoning samples from input queries through the original reasoning model.
    These samples are self-distilled to manifest adaptive reasoning. Finally, root-token policy optimization (\Ours) is applied to the model, refining the policy by updating only the initial branching token (i.e., root-token) to improve difficulty-aware adaptive reasoning.
    }
    \label{fig:landing}
\end{figure}

\begin{figure*}[t]
    \centering
    \includegraphics[width=0.98\linewidth]{ACL2026/figures/1_overview_fig.pdf}
    \caption{Detailed overview of our training framework. After supervised finetuning (SFT) with self-generated correct samples, \Ours~involves updates only the initial decision token with group normalized rewards and advantages.}
    \label{fig:overview}
\end{figure*}

Despite the potential of LRMs, there remains a significant limitation that general LRMs tend to apply this exhaustive explicit reasoning uniformly across all queries, regardless of their intrinsic complexity.
This results in the \textit{overthinking} problem \cite{feng2025efficient, sui2025stop} where substantial computational resources and latency are wasted on trivial questions. Consequently, a surge of emerging research has been driven into striking a balance between reasoning quality and computational efficiency \cite{Zhang2025AdaptThink, Bao2025Thinkless, wu2025arm, wang2025adareasoner}.
While proprietary services such as GPT5 \cite{openai2025gpt5} and Claude \cite{Claude3_5SonnetSystemCard} rely on black-box routing to mitigate overthinking, a novel model-centric alternative is required to internalize adaptive reasoning via policy optimization with open dataset for the open-source community.

While adaptive reasoning models that learn when to reason have been explored via policy optimization, these studies have faced the following challenges:

\begin{enumerate}[leftmargin=12pt, itemsep=2pt]
    \item Training inefficiency. Optimizing over long context trajectories with policy optimization incurs prohibitive VRAM and temporal costs, hindering accessible and iterative development.

    \item Data curation. There is a lack of standardized methods for \textit{baking} difficulty-aware datasets that effectively balance reasoning traces with direct responses across heterogeneous domains.

    \item Initialization gap. Prior works train from scratch or adapt base models, often avoiding adapting the original LRM due to its training sensitivity.
    

\end{enumerate}

To alleviate these challenges, we propose \Ours~(\underline{R}oot-token \underline{P}olicy \underline{O}ptimization), an efficient training method that enables adaptive reasoning by learning only the initial root token (e.g., the decision to open or close the \texttt{<think>} tag). Our method involves two stages (\Autoref{fig:landing}): (1) offline synthesis which generates paired reasoning and non-reasoning responses via forced chat templates, and (2) \Ours, which restricts gradient updates to the pivotal branching decision. 
Our design is complementary to recent token-selective RL that updates only a small subset of decision tokens (e.g., high-entropy forking tokens) rather than the full trajectory~\cite{wang2025beyond8020}; \Ours~pushes this idea to the extreme by optimizing only the root routing token (\Autoref{fig:overview}). 
We demonstrate that \Ours~drastically enhances memory and compute efficiency during training while preserving the model's original capability. 

Our contributions are as follows:

\begin{itemize}[leftmargin=12pt, itemsep=2pt]
    \item We propose \Ours, a novel framework, termed as root-token optimization to enable adaptive reasoning by optimizing only the initial root-token whether to reason, significantly reducing training overhead as well as efficient difficulty-aware routing (\S\ref{sec:method}).

    \item Our empirical results demonstrate \Ours~enables LRMs learning when to reason over diverse LRM families and scales. We also provide an analysis on calibration for an explicit reasoning knob (e.g., target Think\% or routing threshold for root-token probability) (\S\ref{sec:exp}, \S\ref{sec:Analysis}).

    \item We demonstrate that difficulty-aware policies in mathematical domains elicit adaptive reasoning in broader scientific and code domain (\S\ref{sec:Analysis}).
\end{itemize}

